\chapter{An infinitesimal native $r$ with unit $\eta$ and unit $e$}
\label{component:p3763434080c2}
\noindent\textbf{Component }\path{p3763434080c2}\par
{\small\noindent Source: \path{R13/baseline/baseline/baseline/baseline/advances/c1/centered_a_unit_e.tex}\par}
{\small\noindent SHA256: \path{3763434080c27994309135baaeef0216f6916b0507db82e8a646169a209c2016}\par}
{\small\noindent Inventory use: CA1--CA8, complete raw scalar/normal dictionary | CA1-CA8 | C1\_\allowbreak{}BARE\_\allowbreak{}OTHER\_\allowbreak{}TERMINALS; precise subordinate uses retained below\par}
\medskip
\section[An infinitesimal native r with unit eta and unit e]{An infinitesimal native $r$ with unit $\eta$ and unit $e$}

\begin{theorem}[The centered-$A$ unit-$e$ slice]
Use the physical native C1 source with its unchanged original events,
$c=2$, $\lambda=n=0$, all positive openings $a,t,M,J$, and all original
nonnegative signs, including literal $\beta=0$.
Suppose $k,\zeta,r=1-a\eta$ are positive infinitesimals, while
$\eta$ and $e=k\sqrt{\eta/\zeta}$ are positive bounded units.
Then the full parent is tensor SOS at every independent projective test
and every real or complex state. No bound on $k^2/r$ is imposed.
\end{theorem}

The theorem is restricted to unit $e$. Its proof uses $\delta^{-1}$
as a bounded pure coefficient number; it does not assert the cases of
infinitesimal or unbounded $e$.

\paragraph{Centering the sole large event.}
The exact native contact has $Y_c=r u_cW_c$. Consequently
\[
 L_A=Y-r uW
\]
vanishes at contact with no asymptotic replacement. The native formulas
give $a$ a positive unit, $A$ unbounded and $Ar$ a positive unit of
residue $2$. Put $A_*=1/a$, a positive bounded pure coefficient unit.
The complete packet
\[
             4(A-A_*)D_uD_v L_A^2                         \tag{CA1}
\]
is positive and has zero value and full first derivative at the exact
native contact. Retain it as a square.

It is essential to center this packet before taking residues.
The remaining $A$ channel is exactly
\[
 4A_*D_uD_vL_A^2+
 4(Ar)D_uD_v\{2uWY-r u^2W^2\}.                         \tag{CA2}
\]
Only $Ar$ and $r$ occur in its coefficients. In particular this is a
bounded expression for every positive rate of $r$ relative to $k$.

\paragraph{A source chart with no hidden $k^2/r$ quotient.}
Keep the exact contact variables $\delta,k,E=b^2$ from (SD6), and put
\[
 K=1+k^2,\quad \ell=K-aE=Qr>0,\quad
 Q=K^2-\ell k^2,\quad a=(K-\ell)/E.
\]
Here $\delta,E,a$ are positive bounded units and $k,\ell$ are positive
infinitesimals. The exact source substitution is
\[
 \eta=KE/Q,\quad r=\ell/Q,\quad
 \Delta=\delta Q/K,\quad \zeta=\delta^2k^2QE/K,
 \quad R=K^2/Q,\quad T=K-k^2t\zeta.                    \tag{CA3}
\]
All denominators in (CA3) are positive pure coefficient units.
For $i=2,b,3,5,\alpha,d$ let $\widetilde c_i=c_i/\ell$,
defined by cancelling the explicit $L_c=\ell/Q$ factor in the short
native formulas. This cancellation gives rational expressions with
unit denominators; it does not leave an inverse of $\ell$.
For example, writing $H_c=r+\eta$ and $V_c=k^2\eta\zeta$,
\[
 \widetilde c_3=\frac{2KH_cQ_c(1-k^2\eta+tV_c)}Q,
 \qquad \widetilde c_d=\frac{4Q_cP_d}Q.
\]
Their residues are $2E$ for $i=2,b,3,5,\alpha$, and $4$ for $i=d$;
$c_4$ has residue $2$.

Define the following bounded pure coefficient quantities:
\[
 \begin{array}{ll}
 B_2=\dfrac{8a\widetilde c_2}{R\widetilde c_3\Delta^2},&
 B_b=\dfrac{8a\widetilde c_b}{R\widetilde c_3\Delta^2},\\[4pt]
 B_3=8a/T,&B_5=\dfrac{8a\widetilde c_5}{T\widetilde c_3},\\[4pt]
 L=Ar=\dfrac{2a\eta^2c_4}{Q\widetilde c_3},&
 d=\dfrac{2a\eta\widetilde c_d}{\widetilde c_3},\\[4pt]
 \alpha=\dfrac{2a\Delta^2k^6\eta^2\widetilde c_\alpha}
                    {\widetilde c_3}.&
 \end{array}                                                     \tag{CA4}
\]
These formulas explicitly account for the cancellation needed below.

\paragraph{The exact full remainder, before any collection.}
Retain the four original (SD3) packets, divided by the positive source
factor $\Delta^2\Lambda$, and retain (CA1). Use the unchanged pure
testing coordinate $w=(W+uY)/D_u$ and the hard normals $N,M_B$.
For $Z=1$ put $p=t/R$, $h=a/T$ and
\[
 q_A=vD_uw,\quad q_B=v(D_u-uN-puD_uv w),
\]
\[
 W=w+u(M_B+hq_B)/D_u,\qquad
 Y=uw-(M_B+hq_B)/D_u.
\]
The complete remainder is exactly
\begin{align*}
\mathscr G={}&B_2\{N^2-2q_AN+(1-p)q_A^2\}
 +B_bD_v\{N^2+2q_AN+pq_A^2\}\\
&+B_3D_v\{M_B^2+2q_BM_B+hq_B^2\}
 +B_5\{M_B^2-2q_BM_B+(1-h)q_B^2\}\\
&+4A_*D_uD_v(Y-r uW)^2
 +4LD_uD_v(2uWY-r u^2W^2)\\
&+4\alpha D_u(vY-1)^2
 +2dD_u\{(vY-1)^2-D_v\}.                         \tag{CA5}
\end{align*}
There are no normal terms above degree two. Every coefficient of
(CA5) is bounded. The only testing denominator is a power of the
positive unit $D_u$; every source denominator in (CA3)--(CA5) is a
positive unit. The possibly unbounded opening ratios remain exclusively
in the four retained positive packets.

\paragraph{The exact low-weight collection.}
The outer/state lemma in \path{centered_a_outer.tex} reduces the
independent tests to $Z=1$ and $u,w,v$ infinitesimal, with $w>0$ on
the negative pure-core branch. Use the unchanged circuit monomials
\[
 A_0=wv,\quad B_0=\delta v,\quad C_0=\delta Euw,
 \quad D_0=\delta^2k^3E,
\]
\[
 x=B_0-C_0,\quad y=A_0-kC_0,\quad
 z=uB_0-kC_0,\quad h_0=kC_0-D_0.
\]
Let $q_1,q_2$ be the same two separated radical pairings as in the
contact--conductor proof, and let
\[
 T_4=y^2+z^2+h_0^2-2yz+2q_1^2.
\]
Give $(\delta,k,E,\ell,t,u,w,v)$ weights $(0,1,0,1,0,1,1,2)$.
Then $x$ has weight two and $y,z,h_0,q_1,q_2$ have weight three.
Denote the scalar of (CA5) by $S$, and its two complete first normals
by $L_N,L_M$. Direct multiplication of (CA3)--(CA5) through the
indicated weights gives
\[
 \begin{aligned}
 S_{\le6}&=\frac{4P}{E^4\delta^2}x^2
 -\frac{16(k+u)}{E^2\delta^2}xz
 +\frac8{E\delta^2}T_4,\\
 (L_N)_{\le3}&=0,\\
 (L_M)_{\le3}&=
 \frac{8\{E^2+2E-2\ell(E+1)\}}{E^2\delta}\,x,
 \end{aligned}                                                   \tag{CA6}
\]
where the short coefficient polynomial is
\[
 \begin{split}
 P={}&E^2(E+4)-E(E^2+6E+4)\ell+2(E+1)(E+2)\ell^2\\
 &+E^2(8-3E)k^2+4E^2ku+E^2(E+4)u^2.
 \end{split}
\]
The full hard-normal block, from the quadratic part of (CA5), has
residue
\[
       \operatorname{diag}\left(\frac{16}{E\delta^2},
                                \frac{4(E^2+4)}E\right).          \tag{CA7}
\]
Here and in the residue Grams, $E,\delta$ denote their positive
residues. Formula (CA6), rather than a large expansion, is the finite
coefficient identity used by the proof. Its 37 raw scalar monomials
collect to the displayed three terms. The exact checker reconstructs
the six cancelled native coquotients, the entire expression (CA5),
and both complete normals before making this collection.

\paragraph{Why all higher terms have harmless coefficients.}
The scalar of (CA5), after clearing the stated units, belongs to
$(A_0,B_0,C_0)^2+(D_0^2)$. This follows directly from $q_A=wvD_u$,
$q_B|_{N=M_B=0}\in(v)$, and $Y,W$ in (CA5): the centered scalar
$Y-r uW$ lies in $(v,uw)$; its square is in the same raw monomial
square ideal since $\delta,E$ are units. The sole constant state term
comes from $\alpha$, which has the factor $k^6$, and is a bounded
unit multiple of $D_0^2$. Both full first normals belong to
$(A_0,B_0,C_0)$ by the same displayed expression.

All retained packets vanish to first order at the exact contact
\[ (u,w,v,N,M_B)=(k,\delta k,\delta k^2E,0,0). \]
Thus the scalar belongs to the square of its contact ideal and the
first normals belong to that ideal. The unchanged monic
contact--conductor identities therefore apply, over the coefficient
ring in which only the displayed positive units have been inverted.
They never divide by $k$, $\ell$, $M$, $J$, $t$, or a contact difference;
the displayed $1/a$ is the inverse of a known positive unit.
The contact substitutions preserve the new weights. After (CA6) is
subtracted, scalar weights are at least seven and normal weights at
least four. Their lifted coefficients consequently have positive
weights, hence zero residues. Clearing the finitely many unit
denominators before this polynomial lift and restoring them afterward
changes no residue sign and introduces no mixed inverse.

\paragraph{Strictness, full normals, and restoration.}
The leading scalar is
\[
 \frac{4(E+4)}{E^2\delta^2}x^2+\frac8{E\delta^2}T_4.
\]
The standard two-pairing circuit supplies a strict Gram on
$(y,z,h_0,q_1,q_2)$. The only nonzero first-normal residue is
$8(E+2)x/(E\delta)$, the coefficient of the full normal $M_B$.
Accordingly the coupled $(M_B,x)$ Gram is
\[
 \begin{pmatrix}
 4(E^2+4)/E&4(E+2)/(E\delta)\\
 4(E+2)/(E\delta)&4(E+4)/(E^2\delta^2)
 \end{pmatrix},\qquad
 \det=\frac{256}{E^3\delta^2}>0.                      \tag{CA8}
\]
Together with the independent $N$ pivot in (CA7), this gives a
bounded strict eight-root Gram for the entire remainder (CA5).
Restore (CA1), all four original positive packets, and all separated
normalization factors. The outer/state lemma exhausts the remaining
tests; real and imaginary state evaluations give the complex case.

This proves precisely the stated unit-$e$, unit-$\eta$,
infinitesimal-$r$ slice. Combined with the bounded-unit-$r$ theorem,
it covers every positive $r$ when $\eta,e$ are bounded units and
$k,\zeta$ infinitesimal. It makes no assertion about unbounded
$\eta$ or either nonunit alternative for $e$.
